\section{시계 변환과 시간 완결 계약}
\label{sec:time}

시간 계층은 각 원천의 대표 측정시각을 기준 시계로 옮기고, 가용 시각에 따라
창의 완결을 판정한다. LSL은 송신 호스트 시계와 호스트 사이의 offset 관측을 제공하며,
BrainFlow는 장치별 timestamp 의미와 preset별 표본률을 제공한다
\citep{lsl_time,brainflow_data}. 어댑터는 이 관측의 anchor, 시계 영역과 변환 버전을
시간 레코드에 연결한다.

\subsection{유효 구간을 가진 시계 사상}

원천 시계 $\kappa$의 읽기값을 $u$, 기준 시계를 $\kappa_0$로 둔다.
시계 세대 안의 변환 버전 $v$는 유효 구간 $J_v$와 중심 $u_0$를 가진다.
그 구간에서 실제 사상과 추정 사상을 다음과 같이 표현한다.
\begin{equation}
 \begin{aligned}
 \phi_v(u)&=b_v+a_v(u-u_0)+r_v(u),\\
 \widehat\phi_v(u)&=\widehat b_v+\widehat a_v(u-u_0),\qquad u\in J_v.
 \end{aligned}
 \label{eq:clock-affine}
\end{equation}
계약은 실제 사상의 증가성, $\widehat a_v>\delta_a\ge0$과 계수ㆍ잔차의 상한
$|b_v-\widehat b_v|\le\delta_b$,
$|a_v-\widehat a_v|\le\delta_a$,
$|r_v(u)|\le\delta_r$를 포함한다. 따라서 시간 구간은
\begin{equation}
 \begin{aligned}
 \epsilon_v(u)&=\delta_b+|u-u_0|\delta_a+\delta_r,\\
 \Phi_v(u)&=[\widehat\phi_v(u)-\epsilon_v(u),
                  \widehat\phi_v(u)+\epsilon_v(u)]
 \end{aligned}
 \label{eq:clock-interval}
\end{equation}
로 계산한다. 원천 timestamp의 양자화와 jitter를 구간 $U$로 표현하면,
다음 구간 확장을 사용한다.
\begin{equation}
 \Phi_v(U)=\operatorname{hull}\!\bigcup_{u\in U}\Phi_v(u).
 \label{eq:clock-interval-extension}
\end{equation}
변환 레코드는 계수, $J_v$, 가용 시각, clock epoch와 근거를 함께 저장한다.
근거 형식 $\mathsf{Bound}(A,J_v,\epsilon_v)$는 가정 $A$ 아래의 포함 관계를,
$\mathsf{Estimate}(\mathcal D,p,\widehat q_p,n)$는 자료 $\mathcal D$에서 얻은
$p$-분위수 추정과 표본 수를 기술한다.

동일한 affine 구간의 두 읽기값 $u_1,u_2$는 공통 offset을 공유하므로,
상대 시간의 오차는 다음 상한을 갖는다.
\begin{equation}
 \left|
 \phi_v(u_2)-\phi_v(u_1)-\widehat a_v(u_2-u_1)
 \right|
 \le\delta_a|u_2-u_1|+2\delta_r.
 \label{eq:clock-relative}
\end{equation}
두 사상 $\Phi_{12},\Phi_{23}$을 합성할 때 $I_{12}(u)=\Phi_{12}(u)
\subseteq J_{23}$를 요구한다. 중간 구간 전체에 걸친 오차 전파는
\begin{equation}
 \epsilon_{13}(u)\le
 |\widehat a_{23}|\,\epsilon_{12}(u)
 +\sup_{z\in I_{12}(u)}\epsilon_{23}(z)
 \label{eq:clock-compose}
\end{equation}
로 계산한다. 이 규칙은 장치, SDK 호스트와 실행기 시계 사이의 변환 사슬에 적용한다.
Rust의 \code{Instant}는 호스트에서 두 관측 사이의 경과시간을 표현한다.
실행기는 origin과 clock-domain 식별자를 부여하고 외부 시계와의 대응을
식~\eqref{eq:clock-affine}의 레코드로 관리한다\citep{rust_instant}.

\subsection{표본 격자, 패킷과 시간 anchor}

균일 발생 규칙은 원천 시계의 격자
$u_n=u_0+(n-n_0)\Delta$를 사용한다.
affine 구간의 명목 기준 표본률은
$\widehat f_{\mathrm{ref}}=(\widehat a_v\Delta)^{-1}$이다.
스트림 기술자는 채널별 ADC 획득률 $f_{\mathrm{ADC}}$,
장치 방출률 $f_{\mathrm{emit}}$, 어댑터 출력률 $f_{\mathrm{out}}$을 각각 기록한다.
패킷당 채널별 표본 수 $M$이 일정한 경로는
\begin{equation}
 f_{\mathrm{packet}}=\frac{f_{\mathrm{emit}}}{M}
 \label{eq:packet-rate}
\end{equation}
을 만족한다. 패킷 sequence와 식~\eqref{eq:record}의 표본 인덱스 $K_i$는
각자의 출처를 유지한다. 디지털 필터와 재표본화는 출력 격자,
대표시각 anchor와 의존 시간 지지를 함께 갱신한다.

장치의 ADC 관측, SDK 수신, 소프트웨어 명령과 물리 edge는 각각 이름을 가진 anchor이다.
예를 들어 PC 수신 시각을 제공하는 BrainFlow 경로는 해당 수신 경계를
timestamp anchor로 기록한다\citep{brainflow_data}.
동일 pulse의 관측 edge를 대응시키는 경로는 edge 쌍으로 clock mapping을 구성한다.
SpikeGLX의 동기화 경로는 stream별 clock calibration과 공통 sync pulse의
edge 대응을 사용한다\citep{spikeglx_manual}.

\subsection{관측 오차와 대기 자료를 포함한 watermark}

입력 $c$의 원시 레코드에 대해
$I_i^{\mathrm{meas}}=[\ell_i,h_i]
=[\widehat t_i-\epsilon_i,\widehat t_i+\epsilon_i]$로 둔다.
ingress 시각의 관측 구간을
$[\widehat a_i-\eta_i,\widehat a_i+\eta_i]$로 둔다.
한 clock epoch와 시간 뷰에서 지연 계약을 다음과 같이 정의한다.
\begin{equation}
 (\widehat a_i+\eta_i)-(\widehat t_i-\epsilon_i)\le L_c.
 \label{eq:interval-lateness}
\end{equation}
관측 중심시각의 차이가 $\widehat a_i-\widehat t_i\le L_c^0$이고
$\eta_i\le\overline\eta_c$, $\epsilon_i\le\overline\epsilon_c$이면,
$L_c=L_c^0+\overline\eta_c+\overline\epsilon_c$를 사용한다.
따라서 $L_c$는 두 관측의 시계 오차를 함께 포함한다.

현재 기준 시각의 관측 구간을 $s=[s^-,s^+]$로 둔다.
ingress 이후 join 상태에 반영되는 과정에 있는 레코드 집합을
$\operatorname{Pending}_c(s)$라 두고 다음 진행 경계를 계산한다.
\begin{equation}
 \begin{aligned}
 w_c^{\mathrm{in}}(s)&=s^--L_c,\\
 p_c(s)&=\inf\{\ell_i:i\in\operatorname{Pending}_c(s)\},
                  \qquad\inf\varnothing=+\infty,\\
 w_{\Join}(s)&=\min_c\min\{w_c^{\mathrm{in}}(s),p_c(s)\}.
 \end{aligned}
 \label{eq:pending-watermark}
\end{equation}
현재 시각 이후 도착하는 자료는 식~\eqref{eq:interval-lateness}에 따라
$\ell_i>s^--L_c$를 만족한다. 이미 도착한 자료는 $p_c(s)$가 포괄한다.
그러므로 $w_{\Join}$는 join에 앞으로 반영될 자료의 대표시각 하한이다.
구현은 자료, gap, watermark와 epoch 종료 사건의 전달 순서에 이 관계를 유지한다.

창 $W_k=[u_k,v_k)$는 다음 조건에서 완결된다.
\begin{equation}
 \begin{aligned}
 \operatorname{Complete}(W_k,s)&\iff w_{\Join}(s)\ge v_k,\\
 \operatorname{Definite}(i,W_k)&\iff I_i^{\mathrm{meas}}\subseteq W_k,\\
 \operatorname{Possible}(i,W_k)&\iff
                         I_i^{\mathrm{meas}}\cap W_k\ne\varnothing.
 \end{aligned}
 \label{eq:window-complete}
\end{equation}
연산자는 확정ㆍ가능 소속 중 선택한 정책과 경계 모호성 상태를 보관한다.
원시 레코드의 창 할당은 대표 측정 구간을 사용하고,
출력의 전체 시간 지지는 식~\eqref{eq:support}의 조상 합집합을 사용한다.
시간구간 창 출력은 오른쪽 경계 $v_k$를 대표시각 anchor로 선택한다.
표본수 창은 식~\eqref{eq:window}의 마지막 표본 시각을 anchor로 사용한다.
중심 anchor를 사용하는 연산자는 전방 관측 길이 $h_o$를 효과에 기록한다.

ingress 이후 join 반영까지의 상한이 $D_c$이고,
진행 경계 갱신 간격의 상한이 $G$이면 창의 가용 시각은
\begin{equation}
 a_{\mathrm{close}}\le v_k+\max_c(L_c+D_c)+G
 \label{eq:window-availability-bound}
\end{equation}
를 만족한다. 이 식은 기준 시각을 직접 읽는 진행 타이머를 가정한다.
구간 타이머를 사용하는 실행기는 현재 시각의 하한과 실제 시각 사이의 상한을 $G$에 더한다.
가용 시각 인과성은 창 완결 사건을 입력으로 포함하여
식~\eqref{eq:availability-causality}와 연결한다.

\subsection{Marker와 반응시간 계약}

marker는 사건 코드, 발생 원천, 관측 anchor, edge 극성,
측정 구간, 가용 시각과 clock epoch를 가진다.
장치의 디지털 입력 경로는 전압 임계값, 최소 pulse 폭과 표본화 규칙을
하드웨어 계약에 연결한다. 소프트웨어 명령 구간 $I^{\mathrm{cmd}}$와
명령에서 물리 edge까지의 지연 구간 $\Delta_{\mathrm{path}}=[d^-,d^+]$가 주어지면,
명령에 근거한 edge 예측 구간을
\begin{equation}
 I^{\mathrm{edge,pred}}=I^{\mathrm{cmd}}\oplus\Delta_{\mathrm{path}},
 \qquad
 [a,b]\oplus[c,d]=[a+c,b+d]
 \label{eq:marker-delay}
\end{equation}
로 정의한다. edge를 획득한 표본은 그 표본 격자와 clock mapping으로
자신의 $I^{\mathrm{meas}}$를 얻는다.
기준 marker 구간 $[\ell_m,h_m]$, 출력 가용 시각 상한 $\overline a_y$,
반응시간 요구 $D$에 대해
$\overline a_y\le\ell_m+D$를 충분조건으로 사용한다.
실행기는 적용한 지연 구간, 시계 변환과 창 소속 정책을 결과의 계보에 기록한다.
